% !TeX program = xelatex
% !TeX encoding = UTF-8
% HIT 开题报告格式模板：只含占位内容和通用排版示例。
\documentclass[UTF8,fontset=none,scheme=plain,zihao=-4]{ctexart}
\usepackage{hit-proposal}
% 可选：\usepackage[pagenumbers,chinesecaptions]{hit-proposal}
% strictfonts 的使用方法见 README.md；默认无需自行安装中文字体。
% 仅填写封面内容，不在这里修改字号、位置或下划线。
% 公开模板中请保留占位文字；撰写个人报告时使用独立副本。
\proposalsetup{
  reporttype={博士学位开题报告},
  doctype={学位论文},
  title={论文题目（请填写）},
  college={学院名称},
  major={专业学位类别},
  supervisor={导师姓名},
  cosupervisor={（选填）},
  author={研究生姓名},
  studentid={学号},
  date={年\quad 月\quad 日}
}


\begin{document}
\makeproposalcover
\proposaltableofcontents

\section{课题来源及研究的目的和意义}
\subsection{课题来源及研究背景}
在此填写课题来源与研究背景。正文使用统一的首行缩进和行距，不需要手工输入空格来调整段落。

引用格式示例：单篇文献引用\cite{latex-book}，多篇文献引用\cite{latex-book,texbook,ctex-package}。这里的文献仅用于演示排版，正式撰写时应换成实际使用的文献。

\subsection{研究目的及意义}
在此填写研究目的、理论意义或实际应用价值。

\section{国内外研究现状及分析}
\subsection{相关研究现状}
在此按主题介绍已有研究，并给出相应引用。

\subsubsection{三级标题示例}
在此填写这一研究方向的介绍。

\paragraph{四级标题示例}
在此填写细分主题的介绍。本标题用于演示四级编号和目录样式，可按实际需要删除。

\subsection{已有研究的分析}
在此总结已有工作的进展与需要进一步研究的问题。

\section{主要研究内容、实施方案及其可行性论证}
\subsection{主要研究内容}
\begin{enumerate}
  \item 在此填写研究内容一。
  \item 在此填写研究内容二。
  \item 在此填写研究内容三。
\end{enumerate}

\subsection{实施方案}
在此填写研究方法和实施步骤。图~\ref{fig:placeholder} 为图像占位示例；替换图片时建议只指定宽度，以保持原始比例。

\begin{figure}[htbp]
  \centering
  % 替换示例：\includegraphics[width=0.8\linewidth]{figures/example.pdf}
  \fbox{\parbox[c][32mm][c]{0.80\linewidth}{\centering
    图像占位\\[4mm]
    请在此插入图片或技术路线图}}
  \caption{图题示例（请替换）}
  \label{fig:placeholder}
\end{figure}

公式示例见式~\eqref{eq:example}。公式和交叉引用均由 \LaTeX{} 自动编号。
\begin{equation}
  \mathbf{A}\mathbf{x}=\mathbf{b}.
  \label{eq:example}
\end{equation}

\subsection{可行性论证}
在此填写理论、技术、数据或实验条件方面的可行性分析。

\section{预期达到的目标}
在此填写预期目标及其评价方式。

\section{前期的学术研究或者实践工作的结果}
在此填写与本课题相关的前期工作。没有相关内容时，按培养单位要求处理。

\section{进度安排}
表~\ref{tab:schedule} 为三线表和进度安排的格式示例，所有内容均为占位。
\begin{table}[htbp]
  \centering
  \caption{进度安排示例}
  \label{tab:schedule}
  \begin{tabularx}{\linewidth}{@{}p{0.16\linewidth}p{0.20\linewidth}X@{}}
    \toprule
    阶段 & 起止时间 & 工作内容与阶段目标 \\
    \midrule
    第一阶段 & 待填写 & 待填写 \\
    第二阶段 & 待填写 & 待填写 \\
    第三阶段 & 待填写 & 待填写 \\
    \bottomrule
  \end{tabularx}
\end{table}

\section{预期创新性成果}
在此填写预期创新性成果，按实际研究内容展开。

\section{为完成课题已具备和所需的条件和外协计划}
\subsection{已具备的条件}
在此填写已经具备的研究条件。

\subsection{所需条件与外协计划}
在此填写需要补充的条件及合作安排。

\section{研究或者实践工作过程中可能遇到的问题以及解决的措施}
\subsection{可能遇到的问题}
在此填写研究过程中可能遇到的问题。

\subsection{解决措施}
在此填写与上述问题对应的解决措施。

% 本样式自动为参考文献建立编号章节及目录条目。
% 不要再单独添加 \section{参考文献}，也不要编辑生成的 .bbl 文件。
\bibliographystyle{gbt7714-numerical}
\bibliography{references}
\end{document}
